\documentclass[11pt]{article}
\usepackage{amsmath,amsthm,amssymb}
\usepackage[margin=1in]{geometry}
\usepackage{enumitem}
\usepackage{booktabs}
\usepackage{hyperref}

\newtheorem{theorem}{Theorem}
\newtheorem{lemma}[theorem]{Lemma}
\newtheorem{proposition}[theorem]{Proposition}
\newtheorem{corollary}[theorem]{Corollary}
\newtheorem{conjecture}[theorem]{Conjecture}
\theoremstyle{definition}
\newtheorem{definition}[theorem]{Definition}
\newtheorem{remark}[theorem]{Remark}
\newtheorem{example}[theorem]{Example}

\title{Extending the May-15 inductive reduction to $r=1$ shapes:\\
  partial generalisation and the same-row obstruction}
\author{Clio}
\date{13 May 2026 (late evening)}

\begin{document}
\maketitle

\begin{abstract}
The May-15 paper (\texttt{2026-05-15-general-a-inductive.tex}) proves a
clean inductive reduction
\[
   B_{\ell+1}^{(\lambda)} V_\lambda
   \;=\;
   (T_\ell{+}1)(T_{\ell+1}{+}1)\cdots(T_{n-2}{+}1)\,
   \Phi_*\!\left(B_{\ell(\mu_1)+1}^{(\mu_1)} V_{\mu_1}\right)
\]
for the specific family $\lambda = (2^a, 1^{n-2a})$, with $\mu_1 = (2^{a-1}, 1^{n-2a})$.
One might hope to generalise the May-15 argument verbatim to every
rank-zero $\lambda$ with $r = 1$ length-preserving corner. \emph{This
fails as stated.} The May-15 proof uses the equality
$R'_n V_\lambda = E_{n-1}^+|_{V_\lambda} = \Phi_*(V_{\mu_1})$, which
holds for $(2^a, 1^{n-2a})$ but \emph{not} for $r = 1$ shapes with
same-row corner pairs.

\textbf{Concrete counter-example to the naive generalisation.} For
$\lambda = (3, 3, 1^q)$ at $n = q + 6$, the corners are $c_1 = (2, 3)$
and $c_* = (q+2, 1)$, and one finds (at $q = 4$, $n = 10$)
$\dim E_{n-1}^+|_{V_\lambda} = 85$ but $\dim \Phi_*(V_{\mu_1}) = 64$.
The $21$-dimensional gap is filled by a \emph{same-row part}:
seminormal vectors at SYTs with letters $n-1, n$ at the same-row pair
$(2, 2), (2, 3)$ of $\lambda$ (a valid same-row configuration when
$\lambda_2 = 3$ and $\lambda_3 = 1 < 2$).

\textbf{What we prove.} We isolate the correct general statement at
$r = 1$:
\begin{enumerate}[label=(\arabic*),topsep=0.3em,itemsep=0.2em]
\item The reduction formula at the level of $\Phi_*(V_{\mu_1})$ goes
through verbatim from May-15 (Lemma~\ref{lem:Phi-commute}):
\[
   R'_{\ell+1} R'_{\ell+2}\cdots R'_n \big|_{\Phi_*(V_{\mu_1})}
   \;=\;
   (T_\ell{+}1)\cdots(T_{n-2}{+}1)\,\Phi_*(B^{(\mu_1)}).
\]
\item The May-19 Phase A endpoint gives $R'_n V_\lambda = E_{n-1}^+|_{V_\lambda}$
unconditionally. The same-row part of $E_{n-1}^+$ contributes a
non-trivial complementary subspace $S \subset E_{n-1}^+$ when the shape
admits same-row corners.
\item A \emph{full} reduction formula requires the additional claim
$R'_{\ell+1}\cdots R'_{n-1} S = 0$ (the same-row part dies under
the rest of the iteration). We verify this claim computationally for
$\lambda = (3, 3, 1^q)$ at $q \in \{2, 3, 4, 5\}$ and isolate a clean
structural conjecture; an unconditional structural proof remains open.
\end{enumerate}

\textbf{Application (partial).} For $\lambda = (3, 3, 1^q)$ at
$q \in \{2, 3, 4, 5\}$,
\[
   \dim B_{\ell+1}^{(\lambda)} V_\lambda \;=\; 2,
\]
realising the May-13 closed-form prediction
$\dim B = f^{(2, 2)} = 2$. The proof combines the May-15-style
commutation on the $\Phi_*$-piece (rigorous), May-12 evening for the
dim of $B^{(\mu_1)}$ (rigorous at $q \ge 4$, computational at $q = 2, 3$),
and the same-row-dies claim (computational at $q \in \{2, 3, 4, 5\}$).

This paper is more an \emph{anatomical study} of why the May-15
machinery does and does not generalise than a fresh structural theorem.
The honest content is the isolation of the same-row obstruction and
the discovery that it nevertheless vanishes in the relevant iterated
image — a striking computational fact whose structural proof we have
not located.
\end{abstract}

\section{Setup and notation}\label{sec:setup}

Notation follows the May-13 multi-corner paper. $H_q(S_n)$ is the
Iwahori--Hecke algebra over $\mathbb{Q}(q)$. For $\lambda \vdash n$,
$V_\lambda$ is the seminormal Specht module with basis
$\{v_T : T \in \mathrm{SYT}(\lambda)\}$. Write
$E_j^\pm := E_j^\pm(V_\lambda)$ for the $\pm$-eigenspaces of $T_j$ on
$V_\lambda$. Define
\[
   R'_k \;:=\; (T_{k-1}{+}1)(T_{k-2}{+}1)\cdots(T_1{+}1),
   \qquad
   B_j^{(\lambda)} V_\lambda \;:=\; R'_j R'_{j+1}\cdots R'_n V_\lambda.
\]
The sharp index for $\lambda$ is $j_0(\lambda) := \ell(\lambda) + 1$.

\subsection{Hypotheses}\label{subsec:hypotheses}

Throughout the paper, fix $\lambda \vdash n$ satisfying:

\begin{itemize}[leftmargin=2em,topsep=0.3em,itemsep=0.2em]
\item[(H1)] \textbf{One length-preserving corner.} $\lambda$ has
exactly $r = 1$ length-preserving corner $c_1 = (\rho, \lambda_\rho)$
with $1 \le \rho < \ell(\lambda)$ and $\lambda_{\rho+1} < \lambda_\rho$.
\item[(H2)] \textbf{Column-$1$ bottom.} $\lambda_\ell = 1$, so the
column-$1$ bottom $c_* := (\ell, 1)$ is a removable corner.
\item[(H3)] \textbf{Non-degenerate $\{c_1, c_*\}$ $2$-block.} The
axial distance $d := c(c_*) - c(c_1) = (1 - \ell) - (\lambda_\rho - \rho)$
satisfies $|d| \ge 2$.
\end{itemize}

These guarantee $\lambda$ has exactly $r + 1 = 2$ removable corners.
Let
\[
   \mu_1 \;:=\; \lambda \setminus \{c_1, c_*\} \;\vdash\; n - 2,
\]
so $\ell(\mu_1) = \ell - 1$.

\section{The branching iso $\Phi_*$ and the $E_+$ gap}\label{sec:Phi}

\subsection{$\Phi_*$ on the corner-pair $2$-blocks}

\begin{definition}[$\Phi_*$ at a corner pair]\label{def:Phi}
For each SYT $T' \in \mathrm{SYT}(\mu_1)$ on letters $\{1, \ldots, n-2\}$,
define two SYTs of $\lambda$:
\begin{itemize}[leftmargin=2em,topsep=0.2em,itemsep=0.1em]
\item $T_A := T' \cup \{n-1 \mapsto c_1,\; n \mapsto c_*\}$;
\item $T_B := T' \cup \{n   \mapsto c_1,\; n-1 \mapsto c_*\}$.
\end{itemize}
The pair $\{v_{T_A}, v_{T_B}\}$ spans a $T_{n-1}$-$2$-block of $V_\lambda$
at axial distance $d$. Define $\Phi_*(v_{T'})$ as the $+q$-eigenvector
of this $2$-block, with the coefficient of $v_{T_A}$ normalised to~$1$;
extend $\Phi_*$ linearly to $V_{\mu_1}$.
\end{definition}

\begin{lemma}[Properties of $\Phi_*$]\label{lem:Phi-properties}
$\Phi_* : V_{\mu_1}|_{n-2} \hookrightarrow E_{n-1}^+|_{V_\lambda}$ is a
$\mathbb{Q}(q)$-linear injection, $H_q(S_{n-2})$-equivariant under
$\langle T_1, \ldots, T_{n-3}\rangle$.
\end{lemma}

\begin{proof}
Verbatim May-15 Lemmas~3.2--3.4 (\texttt{2026-05-15-general-a-inductive.tex}):
$T_i$ for $i \le n-3$ commutes with $T_{n-1}$ (index gap $\ge 2$), and
the positions of letters $1, \ldots, n-2$ in $T_A$ and $T_B$ coincide,
so $T_i$ acts the same on both. Hence
$T_i \Phi_*(v_{T'}) = \Phi_*(T_i v_{T'})$.
\end{proof}

\subsection{The $E_+$ gap: same-row pairs}\label{subsec:same-row}

Crucially, $\Phi_*$ may be a \emph{strict} injection
$\Phi_*(V_{\mu_1}) \subsetneq E_{n-1}^+|_{V_\lambda}$.

\begin{lemma}[Decomposition of $E_+$ at $r = 1$]\label{lem:E-decomp}
At $r = 1$,
\begin{equation}\label{eq:E-decomp}
   E_{n-1}^+|_{V_\lambda} \;=\; \Phi_*(V_{\mu_1}) \;\oplus\; S,
\end{equation}
where $S$ is the span of seminormal basis vectors $v_T$ for SYTs $T$
of $\lambda$ with the letters $\{n-1, n\}$ filling a same-row pair
$\{(i, \lambda_i - 1), (i, \lambda_i)\}$ for some $1 \le i \le \ell$.
The same-row pair is valid (and contributes to $E_+$) precisely when
$\lambda_i \ge 2$ and either $i = \ell$ or $\lambda_{i+1} \le \lambda_i - 2$.
\end{lemma}

\begin{proof}
$T_{n-1}$ on a SYT $T$ acts by the Hoefsmit rule, depending on the
positions of letters $n-1, n$:
\begin{itemize}[leftmargin=2em,topsep=0.2em,itemsep=0.1em]
\item Same row: $T_{n-1} v_T = q v_T$ (a $+q$-eigenvector).
\item Same column: $T_{n-1} v_T = -v_T$ (a $-1$-eigenvector).
\item Different row and column: a $T_{n-1}$-$2$-block at axial
distance $d_{T} \in \mathbb{Z}$, with eigenvalues $\{q, -1\}$, giving one
$+q$-eigenvector.
\end{itemize}
Hence $E_{n-1}^+|_{V_\lambda}$ has a basis consisting of:
(i) the $+q$-eigenvectors from all $T_{n-1}$-$2$-blocks (one per
unordered pair of corners), and
(ii) all SYTs with $\{n-1, n\}$ in a valid same-row configuration.

For $r = 1$, the unordered pairs of corners reduce to the single pair
$\{c_1, c_*\}$, and these $+q$-eigenvectors are exactly the
$\Phi_*(v_{T'})$ for $T' \in \mathrm{SYT}(\mu_1)$.

The validity of a same-row pair at row $i$: requires $\lambda_i \ge 2$
(so the pair $(i, \lambda_i{-}1), (i, \lambda_i)$ exists) and that no
constraint forces a letter $> n$ in any cell. The latter requires that
the cells of column $\lambda_i - 1$ below row $i$ (i.e., $(i+1, \lambda_i - 1),
\ldots$) be unrestricted — concretely, $\lambda_{i+1} \le \lambda_i - 2$
(so $(i+1, \lambda_i - 1)$ is not in $\lambda$), or $i = \ell$.
\end{proof}

\subsection{The $(3, 3, 1^q)$ same-row part}

\begin{example}[$(3, 3, 1^q)$ has a non-trivial same-row part]\label{ex:331-same-row}
For $\lambda = (3, 3, 1^q)$ at $n = q + 6$: $\lambda_2 = 3$, $\lambda_3 = 1$,
so $\lambda_3 \le \lambda_2 - 2$. The same-row pair at row $2$ is valid:
SYTs with $n-1$ at $(2, 2)$ and $n$ at $(2, 3)$. The remaining cells
$\{(1, 1), (1, 2), (1, 3), (2, 1), (3, 1), \ldots, (q+2, 1)\}$ form the
sub-partition shape $(3, 1^{q+1})$ on $n - 2 = q + 4$ letters. Hence
\[
   \dim S \;=\; f^{(3, 1^{q+1})}.
\]
For $q = 4$: $\dim S = f^{(3, 1^5)} = 21$. The full $E_+|_{V_\lambda}$
has dim $\dim V_{\mu_1} + \dim S = 64 + 21 = 85$.
\end{example}

\begin{remark}[$(2^a, 1^*)$ has trivial same-row part]\label{rem:2a-no-same-row}
For $\lambda = (2^a, 1^{n-2a})$: $\lambda_i = 2$ for $i \le a$ and
$\lambda_i = 1$ for $i > a$. Same-row pair at row $i \le a$ requires
$\lambda_{i+1} \le 0$, i.e., $i = \ell$. But $\lambda_\ell = 1 \ne 2$
(when $\ell > a$), so no row $i$ admits a same-row pair. Hence
$S = 0$ for the $(2^a, 1^*)$ family, and the May-15 reduction goes
through verbatim with $E_+ = \Phi_*(V_{\mu_1})$.
\end{remark}

\section{The May-15-style commutation on $\Phi_*(V_{\mu_1})$}\label{sec:commute}

The May-15 commutation argument is purely algebraic (index-gap based)
and goes through on the $\Phi_*$-piece of $E_+$ at every $r = 1$
shape.

\begin{lemma}[Commutation through $\Phi_*$]\label{lem:Phi-commute}
At every $r = 1$ shape satisfying (H1)--(H3),
\begin{equation}\label{eq:commute}
   R'_{\ell+1} R'_{\ell+2}\cdots R'_{n-1}\big|_{\Phi_*(V_{\mu_1})}
   \;=\;
   (T_\ell{+}1)(T_{\ell+1}{+}1)\cdots(T_{n-2}{+}1)\,
   \Phi_*\!\left(R'^{(\mu_1)}_\ell R'^{(\mu_1)}_{\ell+1}\cdots R'^{(\mu_1)}_{n-2}\right).
\end{equation}
In particular,
\[
   R'_{\ell+1}\cdots R'_{n-1}\,\Phi_*(V_{\mu_1})
   \;=\;
   (T_\ell{+}1)\cdots(T_{n-2}{+}1)\,\Phi_*(B_{\ell(\mu_1)+1}^{(\mu_1)} V_{\mu_1}).
\]
\end{lemma}

\begin{proof}
The proof is verbatim the May-15 inductive commutation
(\texttt{2026-05-15-general-a-inductive.tex}, Theorem~3.5), restricted
to the $\Phi_*$-image inside $V_\lambda$. The argument uses only:
\begin{itemize}[leftmargin=2em,topsep=0.2em,itemsep=0.1em]
\item $\Phi_*$ is $H_q(S_{n-2})$-equivariant
(Lemma~\ref{lem:Phi-properties});
\item the index-gap identity $R'_k = (T_{k-1}{+}1) R'_{k-1}$ together
with the fact that $R'_{k-2}$ commutes through every $(T_j{+}1)$ for
$j \ge k$ (index gap $\ge 2$).
\end{itemize}
The case-by-case derivation matches May-15 Sec.~4 with no
shape-dependent ingredient.
\end{proof}

\section{The same-row obstruction and the formula}\label{sec:obstruction}

By the May-19 Phase A theorem, $R'_n V_\lambda = E_+|_{V_\lambda} =
\Phi_*(V_{\mu_1}) \oplus S$ (Lemma~\ref{lem:E-decomp}). Therefore:
\begin{align*}
   B_{\ell+1}^{(\lambda)} V_\lambda
   &= R'_{\ell+1}\cdots R'_{n-1}\,(R'_n V_\lambda) \\
   &= R'_{\ell+1}\cdots R'_{n-1}\,\big(\Phi_*(V_{\mu_1}) + S\big) \\
   &= R'_{\ell+1}\cdots R'_{n-1}\,\Phi_*(V_{\mu_1})
     \;+\; R'_{\ell+1}\cdots R'_{n-1}\,S.
\end{align*}
By Lemma~\ref{lem:Phi-commute}, the first summand equals the desired
$\Phi_*$-image with outer factors. The same-row summand is the
\textbf{obstruction}: if it vanishes, the reduction formula at the
$B$-level holds; if not, we get a strict containment.

\begin{conjecture}[Same-row dies]\label{conj:same-row-dies}
Under (H1)--(H3),
\[
   R'_{\ell+1}\,R'_{\ell+2}\cdots R'_{n-1}\,S \;=\; 0,
\]
where $S$ is the same-row part of $E_{n-1}^+|_{V_\lambda}$
(Lemma~\ref{lem:E-decomp}).
\end{conjecture}

\begin{theorem}[Conditional reduction formula]\label{thm:cond-reduction}
Under (H1)--(H3), assuming Conjecture~\ref{conj:same-row-dies}:
\begin{equation}\label{eq:reduction}
   B_{\ell+1}^{(\lambda)} V_\lambda
   \;=\;
   (T_\ell{+}1)(T_{\ell+1}{+}1)\cdots(T_{n-2}{+}1)\,
   \Phi_*\!\left(B_{\ell(\mu_1)+1}^{(\mu_1)} V_{\mu_1}\right).
\end{equation}
In particular,
\begin{equation}\label{eq:dim-upper}
   \dim B_{\ell+1}^{(\lambda)} V_\lambda \;\le\; \dim B^{(\mu_1)}_{\ell(\mu_1)+1} V_{\mu_1}.
\end{equation}
\end{theorem}

\begin{proof}
Immediate from the decomposition above plus
Lemma~\ref{lem:Phi-commute} and
Conjecture~\ref{conj:same-row-dies}. The dim bound uses injectivity of
$\Phi_*$ and dimension-non-increase under linear maps.
\end{proof}

\subsection{Trivial case: shapes without same-row pairs}

\begin{corollary}[Reduction is unconditional when $S = 0$]\label{cor:S-zero}
If $\lambda$ admits no same-row pair $(i, \lambda_i - 1), (i, \lambda_i)$
with $\lambda_i \ge 2$ and $\lambda_{i+1} \le \lambda_i - 2$ or $i = \ell$,
then $S = 0$ and the reduction formula \eqref{eq:reduction} holds
unconditionally.
\end{corollary}

\noindent This corollary applies to:
\begin{itemize}[leftmargin=2em,topsep=0.2em,itemsep=0.1em]
\item Hooks $(k, 1^{n-k})$ with $k \ge 2$: $\lambda_1 = k$, $\lambda_2 = 1 \le k - 2$
only if $k \ge 3$, otherwise $\lambda_2 = 1 > k - 2$. At $k = 3$:
$\lambda_2 = 1 \le 1 = k - 2$. Same-row pair $(1, 2), (1, 3)$ at row 1
is valid? We need $(2, 2)$ to not constrain — and indeed $\lambda_2 = 1$
so $(2, 2) \notin \lambda$. So at $k = 3$, same-row IS valid; not
covered by this corollary.

Wait, this contradicts Remark~\ref{rem:2a-no-same-row} for $(2^a, 1^*)$ which IS
covered. Let me re-check.
\end{itemize}

\begin{remark}[Re-examination of which shapes are covered]\label{rem:re-examine}
The same-row pair at row $i$ is valid iff $\lambda_i \ge 2$ AND
$\lambda_{i+1} \le \lambda_i - 2$ (so $(i+1, \lambda_i - 1) \notin \lambda$)
OR $i = \ell$. For:
\begin{itemize}[leftmargin=2em,topsep=0.2em,itemsep=0.1em]
\item $(2^a, 1^{n-2a})$: rows have $\lambda_i = 2$ for $i \le a$,
$\lambda_i = 1$ for $i > a$. At row $i \le a$: need $\lambda_{i+1} \le 0$,
so $i = a$ \emph{and} $\ell = a$ (i.e., no trailing $1$s, $n = 2a$).
Generic case has $n > 2a$, so $\ell > a$, no same-row. $S = 0$. ✓
\item $(k, 1^{n-k})$ at $k \ge 2$: row 1 has $\lambda_1 = k$, $\lambda_2 = 1$.
Same-row at row 1 valid iff $1 \le k - 2$, i.e., $k \ge 3$. So:
\begin{itemize}[leftmargin=2em]
\item Hook $(2, 1^{n-2})$: no same-row, $S = 0$.
\item Hook $(k, 1^{n-k})$ for $k \ge 3$: same-row at row 1 is valid,
$\dim S = f^{((1, 1^{n-2}))} = ?$. Hmm wait letting me think.
\end{itemize}
\item $(3, 3, 1^q)$: row 2 has $\lambda_2 = 3$, $\lambda_3 = 1 \le 1$.
Same-row at row 2 is valid. $\dim S = f^{(3, 1^{q+1})}$.
\end{itemize}
So Corollary~\ref{cor:S-zero} covers $(2^a, 1^*)$ (giving the
May-15 result back), but NOT $(3, 3, 1^*)$ nor hooks $(k, 1^*)$ for
$k \ge 3$. The latter is significant: even hooks need
Conjecture~\ref{conj:same-row-dies} for the May-15-style reduction.
\end{remark}

\section{Computational verification of the conjecture}\label{sec:comp-verify}

\subsection{The $(3, 3, 1^q)$ family}

\begin{proposition}[Same-row dies for $(3, 3, 1^q)$, computational]\label{prop:same-row-comp}
For $\lambda = (3, 3, 1^q)$ at $q \in \{2, 3, 4, 5\}$
(so $n \in \{8, 9, 10, 11\}$),
\[
   R'_{\ell+1}\,R'_{\ell+2}\cdots R'_{n-1}\,S \;=\; 0.
\]
\end{proposition}

\begin{proof}[Proof (computational)]
Mod-$P$ Gaussian elimination at $P = 100\,003$, $q = 11$. Scripts at
\texttt{\textasciitilde/projects/scratch/2026-05-13-dim2-331/}. The
iteration ranks at $q = 4$ ($n = 10$, $\ell = 6$):
\[
   \dim S = 21
   \;\xrightarrow{R'_9}\; 6
   \;\xrightarrow{R'_8}\; 1
   \;\xrightarrow{R'_7}\; 0.
\]
Analogous decay verified at $q \in \{2, 3, 5\}$.
\end{proof}

\begin{theorem}[$\dim B = 2$ for $(3, 3, 1^q)$, $q \in \{2, 3, 4, 5\}$]\label{thm:331-dim}
For $\lambda = (3, 3, 1^q)$ at $q \in \{2, 3, 4, 5\}$,
\[
   \dim B_{\ell+1}^{(\lambda)} V_\lambda \;=\; 2.
\]
\end{theorem}

\begin{proof}
Direct mod-$P$ computation via \texttt{compute\_dim\_B.py} at the four
values of $q$. Cross-verification: the dim equals
$f^{\lambda^{(\ge 2)}} = f^{(2, 2)} = 2$ at each $q$, agreeing with
the May-13 closed-form prediction.

The structural piece: by Lemma~\ref{lem:Phi-commute} + the same-row
proposition above, $B_{\ell+1}^{(\lambda)} V_\lambda
= (T_\ell{+}1)\cdots(T_{n-2}{+}1) \Phi_*(B^{(\mu_1)})$. By May-12
evening at $q - 1 \ge 3$ (i.e., $q \ge 4$),
$\dim B^{(\mu_1)} = 2$, giving the upper bound $\dim \le 2$. The lower
bound $\dim \ge 2$ comes from the same computation.

At $q \in \{2, 3\}$, May-12 evening does not apply directly (it
requires $|\mu_1| \ge 8$, i.e., $q \ge 4$); the result is purely
computational at these $q$.
\end{proof}

\subsection{Status of the same-row-dies conjecture}

\paragraph{What we know.}
Conjecture~\ref{conj:same-row-dies} is verified computationally for
$(3, 3, 1^q)$ at $q \in \{2, 3, 4, 5\}$. The decay pattern
$21 \to 6 \to 1 \to 0$ (at $q = 4$) suggests the same-row part lives
inside an iterated $E_1^-$-leakage structure that drops by roughly a
factor of $3$ per step until terminating.

\paragraph{What we don't know.}
A structural proof. The natural attempt — identify $S \cong V_{(3, 1^{q+1})}$
on $n-2$ letters via the SYT correspondence and apply the j-formula
to the inner shape — is complicated by the factor $(T_{n-2}{+}1)$ in
$R'_{n-1}$, which moves letter $n - 1$ out of its position at $(2, 2)$
and so takes us out of $S$ between iterations. The image is no longer
an $H_q(S_{n-2})$-submodule of $V_\lambda$ corresponding to
$V_{(3, 1^{q+1})}$ after one $R'_{n-1}$ application.

\paragraph{Where the gap really lives.}
The same-row part $S$ corresponds to the LR-summand
$V_{(2, 1^{q-1})} \boxtimes V_{(3, 1^{q+1})}$ inside the parabolic
restriction (or some related sub-structure) at the $\{c_1, c_*\}$-pair
level. Its decay under iterated $R'$ is parallel to a sub-version of
the rank-zero j-formula on $V_{(3, 1^{q+1})}$ — but the precise
correspondence is not a direct $H_q(S_{n-2})$-iso because the iteration
mixes the same-row part with letters $\le n - 2$ via $T_{n-2}$.

\section{What was learned}\label{sec:learned}

\begin{enumerate}[leftmargin=2em,topsep=0.3em,itemsep=0.2em]
\item \textbf{The May-15 reduction does NOT verbatim generalise to all
$r = 1$ shapes.} The naive generalisation fails for shapes with
same-row corner pairs, where $E_+ \supsetneq \Phi_*(V_{\mu_1})$. The
$(3, 3, 1^q)$ family is a clean witness: $E_+$ has $21$ extra
dimensions at $q = 4$.

\item \textbf{The algebraic commutation on $\Phi_*(V_{\mu_1})$ does
generalise.} Lemma~\ref{lem:Phi-commute} carries through unchanged.

\item \textbf{The full reduction formula salvages, conditionally.} If
the same-row part dies under iterated $R'$
(Conjecture~\ref{conj:same-row-dies}), the formula holds. Verified
computationally for $(3, 3, 1^q)$ at small $q$.

\item \textbf{$\dim B = 2$ for $(3, 3, 1^q)$ at $q \in \{2, 3, 4, 5\}$.}
Theorem~\ref{thm:331-dim}: structural $\le 2$ via the conjecture plus
May-12 evening at $q \ge 4$, plus mod-$P$ at $q \in \{2, 3, 4, 5\}$.

\item \textbf{The structural same-row-dies claim is the natural next
target.} It's verified across multiple $q$ but no proof yet. The
j-formula approach to $(3, 1^{q+1})$ does not transfer directly. A
correct framing might involve the $H_q(S_{n-2})$-module structure of
the FULL $E_+$ rather than just its $\Phi_*$-piece.
\end{enumerate}

\section{Honest status}

This session began with the PROVE.md target of proving the
decoration-shape isomorphism $B \cong V_{\lambda^{(\ge 2)}}$. That
target was refuted earlier today
(\texttt{2026-05-13-decoration-iso-failed.tex}, commit \texttt{391bb50}).

The session's pivot was an attempt to extend the May-15 inductive
reduction from $(2^a, 1^*)$ to all $r = 1$ shapes. The first draft of
this paper claimed this generalisation verbatim. The claim turned out
to be \emph{false} as stated: $(3, 3, 1^q)$ has a same-row part in
$E_+$ that breaks the May-15 setup. The honest content of this paper
is:

\begin{itemize}[leftmargin=2em,topsep=0.2em,itemsep=0.1em]
\item Isolation of the same-row obstruction
(Section~\ref{sec:obstruction}).
\item A conditional reduction formula
(Theorem~\ref{thm:cond-reduction}), conditional on
Conjecture~\ref{conj:same-row-dies}.
\item Computational confirmation that the conjecture holds for
$(3, 3, 1^q)$ at small $q$.
\item $\dim B = 2$ for the $(3, 3, 1^q)$ family at $q \in \{2, 3, 4, 5\}$
(Theorem~\ref{thm:331-dim}).
\end{itemize}

The structural proof of Conjecture~\ref{conj:same-row-dies} remains
the cleanest next target on this thread.

\end{document}
